Modern data systems process enormous amount of data and selecting the appropriate data structure can have major implications in the responsiveness, availability and scalability of such systems. 
Traditional systems have certain limitations and thus more advanced Key-Value systems are used in most applications today. 
During updates, inserts and deletes, LSM trees utilize the main memory of the system to achieve better performance. 
Values are not deleted or updated in place. 
Each delete requires a new entry refered as a 'tombstone' that marks all previous entries with the same key as deleted.
In our knowledge, none of the major Key-Value store developers have implemented range deletes except RocksDB.
Their implementation use the so called 'range tombstones' which alltogether consist a data-structure that can answer wheather a requested key is deleted by a range delete or not. 
For the database to construct such a response multiple I/Os are needed.
Our method tries to reduce these I/Os by efficiently storing past results in memory in an opportunistic way, thus increasing responsivnes, since I/Os majorly affect performance.
